
\pagestyle{empty}\thispagestyle{empty}
\vspace*{\fill}
\begin{center}
\fontsize{52}{52}\selectfont
\textsc{scratch au collège}

  \hfil
  
  \includegraphics[width = 6cm]{logoScratchAuCollege}
\end{center}
\vfill
\begin{center}
\huge
\textsc{algorithmes \  et \  programmation}
\end{center}
\begin{center}
\LogoExoSept{2}
\end{center}
%\clearemptydoublepage
\clearpage

\thispagestyle{empty}

\vspace*{\fill}
\section*{À la découverte du code}

Scratch est un logiciel idéal pour apprendre à programmer. Il a été spécialement conçu pour les enfants et les débutants. La programmation avec Scratch est ludique car il est facile de faire de beaux dessins et des petits jeux. En plus, la programmation est simple, car il suffit de déplacer des blocs pour écrire son code.

\bigskip 

Pourquoi apprendre à coder ? Pour utiliser un ordinateur, je n'ai pas besoin de savoir le programmer ! C'est comme pour les voitures, je n'ai pas besoin de connaître la mécanique pour conduire. Mais, dans le monde qui nous entoure, l'informatique est partout, dans les ordinateurs bien sûr, mais aussi dans nos téléphones, et en fait dans tous les appareils électroniques, et bientôt, ce seront les ordinateurs qui piloteront les voitures ! Il est donc indispensable d'apprendre à parler le langage des ordinateurs. 

\bigskip

Les langages pour programmer un ordinateur sont nombreux, mais une fois qu'un langage est bien compris, les autres s'apprennent plus vite.
Scratch est facile à prendre en main et il permet d'aborder bon nombre de situations de programmation. %de faire tout ce dont on a besoin.
Avec Scratch, la programmation devient un jeu et votre ordinateur un compagnon.
Alors, prêts à programmer ?

\vspace*{\fill}

%\newpage
\thispagestyle{empty}
\addtocontents{toc}{\protect\setcounter{tocdepth}{1}}
\tableofcontents
